\chapter*{Özet}
\addcontentsline{toc}{chapter}{Özet}

\begin{center}
\textbf{TOPLU TAŞIMADA OTOBÜS VARIŞ SÜRELERİNİN TAHMİNİ}
\end{center}

\vspace{0.5cm}

%% Edit below this line
Bu çalışma, İzmir ESHOT toplu taşıma ağında duraklar arası otobüs seyahat sürelerini tahmin etmek için bağlam duyarlı bir makine öğrenmesi sistemi sunmaktadır. Şehrin açık ulaşım Uygulama Programlama Arayüzünden (API) gerçek zamanlı otobüs konumlarını, statik GTFS (General Transit Feed Specification) tarife verisini ve bağlamsal hava durumu ile trafik bilgilerini toplayan, uçtan uca ve tamamen tekrarlanabilir bir işlem hattı geliştirilmiştir. Bu kaynaklardan temiz duraktan-durağa seyahat süresi segmentleri çıkarılmış ve zamansal, mekânsal, tarifeye dayalı, tarihsel ve duraklama süresi (dwell time) özellikleriyle zenginleştirilmiştir. Veri seti, üç otobüs hattı üzerinde yaklaşık üç aylık bir sürede bağımsız olarak toplanmıştır.

Ortak ve kronolojik olarak bölünmüş bir test kümesi üzerinde çeşitli modeller eğitilip karşılaştırılmıştır: naif ve tarihsel temel modeller, Doğrusal Regresyon, Rastgele Orman (Random Forest), XGBoost ve çift girdili bir LSTM (Long Short-Term Memory) ağı. En iyi modeller, tarife temelli baz çizgisinin çok altında, yaklaşık $0{,}43$ dakikalık bir Ortalama Mutlak Hata elde etmektedir. Titiz bir istatistiksel analiz, klasik XGBoost modeli ile derin LSTM'in istatistiksel olarak ayırt edilemez olduğunu, her ikisinin de Rastgele Orman'dan anlamlı ölçüde daha iyi olduğunu göstermektedir; bu nedenle XGBoost en pratik tekil model olarak belirlenmiştir. Yöntem ayrıca üç hatta da tutarlı şekilde genellenmektedir.

Proje, önerilen her iyileştirmenin yalnızca ölçülebilir bir gelişme sağladığında korunduğu A/B temelli bir geliştirme süreciyle metodolojik dürüstlüğe güçlü bir vurgu yapmaktadır. Bu süreç üç temel bulgu üretmiştir: ulaşılabilir doğruluk, 30 saniyelik örnekleme aralığının dayattığı bir ölçüm tabanıyla sınırlıdır; model karmaşıklığını azaltmak doğruluğu düşürmemektedir; ve adil bir aynı-test-kümesi karşılaştırması, derin modelin en iyi olduğuna dair önceki izlenimi düzeltmektedir. GTFS tarifesinden ve GPS'ten türetilen duraklama süresi özelliklerinin en bilgilendirici girdiler arasında olduğu doğrulanmıştır.

%% Until here
\vfill
%% Edit after {Anahtar Kelimeler:}
\textbf{Anahtar Kelimeler:} Otobüs Varış Süresi Tahmini, Makine Öğrenmesi, XGBoost, LSTM, GTFS, Toplu Taşıma, İzmir.
\clearpage
